% Options for packages loaded elsewhere
\PassOptionsToPackage{unicode}{hyperref}
\PassOptionsToPackage{hyphens}{url}
\PassOptionsToPackage{dvipsnames,svgnames,x11names}{xcolor}
%
\documentclass[
  letterpaper,
  DIV=11,
  numbers=noendperiod]{scrartcl}

\usepackage{amsmath,amssymb}
\usepackage{iftex}
\ifPDFTeX
  \usepackage[T1]{fontenc}
  \usepackage[utf8]{inputenc}
  \usepackage{textcomp} % provide euro and other symbols
\else % if luatex or xetex
  \usepackage{unicode-math}
  \defaultfontfeatures{Scale=MatchLowercase}
  \defaultfontfeatures[\rmfamily]{Ligatures=TeX,Scale=1}
\fi
\usepackage{lmodern}
\ifPDFTeX\else  
    % xetex/luatex font selection
\fi
% Use upquote if available, for straight quotes in verbatim environments
\IfFileExists{upquote.sty}{\usepackage{upquote}}{}
\IfFileExists{microtype.sty}{% use microtype if available
  \usepackage[]{microtype}
  \UseMicrotypeSet[protrusion]{basicmath} % disable protrusion for tt fonts
}{}
\makeatletter
\@ifundefined{KOMAClassName}{% if non-KOMA class
  \IfFileExists{parskip.sty}{%
    \usepackage{parskip}
  }{% else
    \setlength{\parindent}{0pt}
    \setlength{\parskip}{6pt plus 2pt minus 1pt}}
}{% if KOMA class
  \KOMAoptions{parskip=half}}
\makeatother
\usepackage{xcolor}
\setlength{\emergencystretch}{3em} % prevent overfull lines
\setcounter{secnumdepth}{-\maxdimen} % remove section numbering
% Make \paragraph and \subparagraph free-standing
\makeatletter
\ifx\paragraph\undefined\else
  \let\oldparagraph\paragraph
  \renewcommand{\paragraph}{
    \@ifstar
      \xxxParagraphStar
      \xxxParagraphNoStar
  }
  \newcommand{\xxxParagraphStar}[1]{\oldparagraph*{#1}\mbox{}}
  \newcommand{\xxxParagraphNoStar}[1]{\oldparagraph{#1}\mbox{}}
\fi
\ifx\subparagraph\undefined\else
  \let\oldsubparagraph\subparagraph
  \renewcommand{\subparagraph}{
    \@ifstar
      \xxxSubParagraphStar
      \xxxSubParagraphNoStar
  }
  \newcommand{\xxxSubParagraphStar}[1]{\oldsubparagraph*{#1}\mbox{}}
  \newcommand{\xxxSubParagraphNoStar}[1]{\oldsubparagraph{#1}\mbox{}}
\fi
\makeatother


\providecommand{\tightlist}{%
  \setlength{\itemsep}{0pt}\setlength{\parskip}{0pt}}\usepackage{longtable,booktabs,array}
\usepackage{calc} % for calculating minipage widths
% Correct order of tables after \paragraph or \subparagraph
\usepackage{etoolbox}
\makeatletter
\patchcmd\longtable{\par}{\if@noskipsec\mbox{}\fi\par}{}{}
\makeatother
% Allow footnotes in longtable head/foot
\IfFileExists{footnotehyper.sty}{\usepackage{footnotehyper}}{\usepackage{footnote}}
\makesavenoteenv{longtable}
\usepackage{graphicx}
\makeatletter
\newsavebox\pandoc@box
\newcommand*\pandocbounded[1]{% scales image to fit in text height/width
  \sbox\pandoc@box{#1}%
  \Gscale@div\@tempa{\textheight}{\dimexpr\ht\pandoc@box+\dp\pandoc@box\relax}%
  \Gscale@div\@tempb{\linewidth}{\wd\pandoc@box}%
  \ifdim\@tempb\p@<\@tempa\p@\let\@tempa\@tempb\fi% select the smaller of both
  \ifdim\@tempa\p@<\p@\scalebox{\@tempa}{\usebox\pandoc@box}%
  \else\usebox{\pandoc@box}%
  \fi%
}
% Set default figure placement to htbp
\def\fps@figure{htbp}
\makeatother

\KOMAoption{captions}{tableheading}
\makeatletter
\@ifpackageloaded{caption}{}{\usepackage{caption}}
\AtBeginDocument{%
\ifdefined\contentsname
  \renewcommand*\contentsname{Tabla de contenidos}
\else
  \newcommand\contentsname{Tabla de contenidos}
\fi
\ifdefined\listfigurename
  \renewcommand*\listfigurename{Listado de Figuras}
\else
  \newcommand\listfigurename{Listado de Figuras}
\fi
\ifdefined\listtablename
  \renewcommand*\listtablename{Listado de Tablas}
\else
  \newcommand\listtablename{Listado de Tablas}
\fi
\ifdefined\figurename
  \renewcommand*\figurename{Figura}
\else
  \newcommand\figurename{Figura}
\fi
\ifdefined\tablename
  \renewcommand*\tablename{Tabla}
\else
  \newcommand\tablename{Tabla}
\fi
}
\@ifpackageloaded{float}{}{\usepackage{float}}
\floatstyle{ruled}
\@ifundefined{c@chapter}{\newfloat{codelisting}{h}{lop}}{\newfloat{codelisting}{h}{lop}[chapter]}
\floatname{codelisting}{Listado}
\newcommand*\listoflistings{\listof{codelisting}{Listado de Listados}}
\makeatother
\makeatletter
\makeatother
\makeatletter
\@ifpackageloaded{caption}{}{\usepackage{caption}}
\@ifpackageloaded{subcaption}{}{\usepackage{subcaption}}
\makeatother

\ifLuaTeX
\usepackage[bidi=basic]{babel}
\else
\usepackage[bidi=default]{babel}
\fi
\babelprovide[main,import]{spanish}
% get rid of language-specific shorthands (see #6817):
\let\LanguageShortHands\languageshorthands
\def\languageshorthands#1{}
\usepackage{bookmark}

\IfFileExists{xurl.sty}{\usepackage{xurl}}{} % add URL line breaks if available
\urlstyle{same} % disable monospaced font for URLs
\hypersetup{
  pdftitle={Optimización Agricola},
  pdfauthor={Ian Contreras; Elvis Lagrange},
  pdflang={es},
  colorlinks=true,
  linkcolor={blue},
  filecolor={Maroon},
  citecolor={Blue},
  urlcolor={Blue},
  pdfcreator={LaTeX via pandoc}}


\title{Optimización Agricola}
\author{Ian Contreras \and Elvis Lagrange}
\date{2025-01-11}

\begin{document}
\maketitle

\renewcommand*\contentsname{Tabla de contenidos}
{
\hypersetup{linkcolor=}
\setcounter{tocdepth}{3}
\tableofcontents
}

\subsection{Planificación de un proyecto agroforestal en la República
Dominicana (Cambiar el nombre de este
inciso)}\label{planificaciuxf3n-de-un-proyecto-agroforestal-en-la-repuxfablica-dominicana-cambiar-el-nombre-de-este-inciso}

El sector agrícola constituye uno de los sectores más importante para la
República, pues en la encargada de proveer de los alimentos esenciales
para los dominicanos, tanto los productores agropecuarios como otros
actores son piezas claves en este proceso. Los productores dominicanos
producen una enorme variedad de rubros agropecuarios lo que lo convierte
en uno de los sectores más diversos.

Gran parte de los poductores agrícolas en la República Dominicana,
cultivan diferentes variedades de rubros agrícolas, con el objetivo de
mitigar los riesgos asociados a los peligros climáticos y a su vez,
garantizar la protección de su producción ante fluctuaciones del
mercado. Esto puede dificultar el manejo óptimo de la producción
(eleeción óptimo de los rubros, manejo de pesticidas, uso de
fertilizantes etc).

Ante esta situación es crucial para los productores agrícolas saber la
combinanción de rubros agrícolas que maximicen sus beneficios. Para ello
se pueden incorporar diferentes métodos matemáticos.En el trabajo de
Miao et al.~(2023), se utilizó el modelo de portafolio para ayudar a los
agricultores a distribuir eficientemente sus recursos (tierra, agua y
trabajo) entre diferentes cultivos, maximizando beneficios y minimizando
riesgos. Este método, inspirado en inversiones financieras, mostró que
asignar más recursos a cultivos como chile y melón mejoraba las
ganancias y la eficiencia, sin aumentar el riesgo. Es una herramienta
útil para que los agricultores optimicen su producción de forma
calculada y sostenible.

Empezar aquí

\section{\textless\textless\textless\textless\textless\textless\textless{}
HEAD}\label{head}

\subsection{Metodología}\label{metodologuxeda}

Para la optimización de portafolios, es necesario conocer los
rendimientos de los activos. En este caso, se emplearán como activos 12
rubros agrícolas: arroz, maíz, papa blanca, batata, yautía blanca, ñame,
berenjena criolla, auyama, aguacate, guineo verde y plátano.

El cálculo de los rendimientos se basa en la productividad promedio
anual de cada cultivo, medida en kilogramos por tarea (kg/tarea), los
precios promedio anuales al productor (en
RD\(/kg) de cada uno de los cultivos, los costos anuales de producción de los rubros agrícolas por tarea y el ciclo de cosecha de cada bien agrícola. Todos estos datos fueron extraídos del Ministerio de Agricultura de la República Dominicana. A continuación, se presenta la fórmula empleada para el cálculo de los rendimientos\)\$
\frac{\text{Productividad} \cdot \text{Precio} - \text{costo}}{\text{costo}}
\cdot (\text{ciclo}) \$\$

El uso de los costos totales como denominador se justifica en su
capacidad para reflejar de manera más precisa los gastos recurrentes
relacionados a la producción, especialmente en ausencia de datos claros
sobre costos fijos iniciales. Además, los costos normales son más
representativos de la variabilidad en la producción y a su vez
simplifcan el análisis al brindar una medida dinámica y facílmente
calculable.

\subsection{Descripción general de los rendimientos de los rubros
agrícolas}\label{descripciuxf3n-general-de-los-rendimientos-de-los-rubros-agruxedcolas}

Al analizar los rendimientos anuales de los rubros agrícolas en el
período 2002-2022, se observa patrones diversos en su comportamiento.
Rubros como el arroz, ñame, plátano y berenjena muestran una tendencia
relativamente estable, con rendimientos que se mantienen
mayoritariamente entre -2\% y 2\%, lo cual evidencia una dinámica de
producción consistente y menos expuesta a fluctuaciones externas
significativas. Por otro lado, cultivos como el aguacate y la auyama
destacan por su mayor volatilidad, con el aguacate registrando un pico
notable que alcanza hasta un 12\% alrededor de 2016, seguido de una
caída pronunciada. La auyama, por su parte, muestra una tendencia
marcadamente alcista en los últimos años, alcanzando rendimientos
superiores al 10\% hacia el final del período, lo que sugiere mejoras
significativas en sus condiciones de producción. La yuca presenta un
comportamiento moderadamente volátil, con fluctuaciones que se mueven
entre 0\% y 2\%, mostrando picos ocasionales cercanos al 2\% a lo largo
del período. En cuanto a la papa, esta exhibe una volatilidad notable
con rendimientos que oscilan entre 0\% y 3\%, mostrando una serie de
picos y valles que reflejan la sensibilidad del cultivo a diversos
factores externos. El maíz muestra una tendencia gradualmente positiva
aunque modesta, mientras que la batata mantiene un comportamiento
relativamente estable hasta un incremento notable hacia 2020, superando
el 2\%. La yautía presenta un patrón particular con un pico
significativo al inicio del período, seguido de una estabilización en
niveles más moderados.

\begin{quote}
\begin{quote}
\begin{quote}
\begin{quote}
\begin{quote}
\begin{quote}
\begin{quote}
dfeeb6831d2bef1a689ceb32428ff0599c588d6e
\end{quote}
\end{quote}
\end{quote}
\end{quote}
\end{quote}
\end{quote}
\end{quote}

\subsection{Teoría de portafolio de Markowitz aplicada a la
agricultura}\label{teoruxeda-de-portafolio-de-markowitz-aplicada-a-la-agricultura}

La Teoría Moderna de Portafolios (MPT), desarrollada por Harry Markowitz
en 1952 (galardonado con el Premio Nobel de Economía en 1990), parte de
la idea de que los agentes ---inversores o productores--- son aversos al
riesgo y buscan la mejor relación entre ganancia esperada y riesgo
asociado. Para ello, la MPT propone cuantificar la volatilidad de cada
activo ---o cultivo, en el contexto agrícola de la presente
investigación--- y examinar las correlaciones entre ellos. La clave
radica en combinar rubros con baja correlación, de modo que el riesgo
total disminuya sin sacrificar rendimientos. En el plano financiero,
este conjunto de combinaciones da lugar a la ``frontera eficiente'', que
representa todos los portafolios que brindan el mayor rendimiento
esperado para cada nivel de riesgo o, de manera equivalente, el menor
riesgo para un nivel de rendimiento específico.

Trasladada al sector agrícola, la metodología de Markowitz permite
determinar la proporción óptima de recursos asignados a cada cultivo,
equilibrando la rentabilidad potencial con la volatilidad inherente a
cada rubro. Entre los portafolios que conforman la frontera eficiente,
el denominado ``portafolio óptimo'' es aquel que maximiza el beneficio
ajustado al riesgo y, por ende, equilibra de manera adecuada la
rentabilidad y la volatilidad esperada. Esta herramienta resulta
especialmente útil para los productores que buscan decisiones informadas
y sostenibles en la asignación de sus recursos. En este punto, resulta
relevante mostrar cómo se implementa este método en la práctica.

La representación matemática de este enfoque se basa en dos elementos
fundamentales: la rentabilidad media del conjunto de cultivos y la
medición del riesgo a través de la varianza. Para el primer caso, se
considera que cada cultivo \(i\) recibe una fracción \(x_i\) de los
recursos (tierra, agua y mano de obra), de modo que
\(\sum_{i=1}^n x_i = 1\). Con estos valores, la rentabilidad media del
``portafolio'' agrícola ---un promedio ponderado de las ganancias
esperadas para toda la producción--- se expresa como:

\[
R = \sum_{i=1}^n x_i r_i,
\]

donde \(r_i\) es el rendimiento estimado del cultivo \(i\), estimado a
partir de la relación entre la productividad promedio anual de cada
cultivo, medida en kilogramos por tarea (kg/tarea), los precios promedio
anuales al productor (en RD\$/kg) de cada uno de los cultivos, los
costos anuales de producción de los rubros agrícolas por tarea y el
ciclo de cosecha de cada bien agrícola.

En este contexto, se considera \(V\) como la medida del riesgo de la
inversión, el cual se expresa mediante la varianza de los retornos.
Siguiendo la notación anterior, se define

\[
V = \sum_{i=1}^n \sum_{j=1}^n x_i x_j \sigma_{ij},
\]

donde \(x_i\) y \(x_j\) representan la proporción de recursos destinados
a los cultivos \(i\) y \(j\), respectivamente, mientras que
\(\sigma_{ij}\) es la covarianza entre los retornos netos de ambos
rubros. El riesgo de la inversión (\(V\)) refleja hasta qué punto las
variaciones de un cultivo influyen en las de otro mayor. Cuando
\(\sigma_{ij}\) es alta, significa que ambos rubros fluctúan de manera
conjunta, lo que aumenta el riesgo total si uno de ellos experimenta un
descenso significativo.

La \textbf{covarianza} entre dos cultivos, \(\sigma_{ij}\), mide el
grado en que los retornos de ambos cultivos varían conjuntamente.
Matemáticamente, puede expresarse como:

\[
\sigma_{ij} = E\left[(r_i - E(r_i))(r_j - E(r_j))\right] = E(r_i r_j) - E(r_i)E(r_j),
\]

donde \(r_i\) y \(r_j\) representan el retorno neto por unidad de
recurso de los cultivos \(i\) y \(j\), respectivamente.

Si los rendimientos de dos cultivos tienden a variar en la misma
dirección, la covarianza entre ellos será positiva; por el contrario, si
sus rendimientos suelen variar en direcciones opuestas, la covarianza
será negativa. Una covarianza negativa facilita una diversificación más
eficiente, ya que permite al productor seleccionar cultivos con
rendimientos poco correlacionados, lo que reduce el riesgo total en la
asignación de recursos

En este enfoque, es necesario plantear los siguientes dos problemas de
optimización:

Para minimizar el riesgo de la producción, bajo el requisito de
garantizar un nivel mínimo de rentabilidad, se resuelve:

\[
\text{Min}(V) = \sum_{i=1}^n \sum_{j=1}^n x_i x_j \sigma_{ij} \quad \text{con} \quad 
\begin{cases}
\sum_{i=1}^n x_i r_i \geq \theta, \\
x_i \geq 0, \\
\sum_{i=1}^n x_i = 1,
\end{cases}
\]

donde \(\theta\) representa el retorno mínimo requerido por el productor
agrícola.

Para maximizar la rentabilidad, bajo el supuesto de que el riesgo no
supere un cierto umbral, se plantea:

\[
\text{Max}(R) = \sum_{i=1}^n x_i r_i \quad \text{sujeto a}:
\]

\begin{enumerate}
\def\labelenumi{\arabic{enumi}.}
\tightlist
\item
  \(\sum_{i=1}^n \sum_{j=1}^n x_i x_j \sigma_{ij} \leq \gamma\),\\
\item
  \(x_i \geq 0\),\\
\item
  \(\sum_{i=1}^n x_i = 1\).
\end{enumerate}

donde \(\gamma\) es el nivel máximo de riesgo que el productor está
dispuesto a asumir.

Al resolver estos problemas de optimización, se obtienen asignaciones de
cultivos que minimizan la volatilidad (varianza) para un determinado
umbral de rentabilidad o, de manera equivalente, que maximizan los
retornos sin sobrepasar el riesgo prefijado. En consecuencia, se estaría
construyendo la frontera eficiente de asignaciones de recursos a los
distintos rubros agrícolas.

\subsection{Justificación
metodológica}\label{justificaciuxf3n-metodoluxf3gica}

Los agricultores suelen adoptar un enfoque empírico y hasta
``heurístico'' al asignar recursos a sus distintos cultivos, basándose
en su experiencia y en un análisis intuitivo de los cultivos que
consideran más rentables. Este enfoque, aunque práctico en el corto
plazo, presenta deficiencias estructurales. Los rendimientos pasados no
garantizan resultados futuros debido a factores climáticos, enfermedades
o fluctuaciones del mercado, lo que genera alta incertidumbre en su
capacidad predictiva. En contraste, la MPT incorpora la varianza y la
covarianza de los rendimientos históricos para estimar la rentabilidad
esperada ajustada al riesgo, permitiendo decisiones más robustas frente
a incertidumbres futuras. Además, mientras la selección intuitiva de
cultivos maximiza beneficios en el corto plazo, ignora los beneficios de
la diversificación, que reduce el riesgo total y mejora la rentabilidad
ajustada en el largo plazo.

La programación lineal (PL), aunque común en la optimización de recursos
agrícolas, también presenta restricciones significativas. Asume una
rentabilidad fija y determinista para cada cultivo, limitando su
capacidad para modelar la incertidumbre inherente a la producción
agrícola. Por el contrario, la MPT optimiza en función de la
rentabilidad esperada ponderada por el riesgo, considerando también las
interacciones entre cultivos a través de las covarianzas. Esto permite
capitalizar los efectos diversificadores de combinar cultivos con baja o
negativa correlación, algo que la PL no puede modelar. Además, la MPT
proporciona un marco más flexible que permite adaptaciones dinámicas a
condiciones cambiantes, mientras que la PL tiende a ofrecer soluciones
rígidas y poco realistas, como la asignación completa de recursos a un
solo cultivo.

Adoptar la MPT en la planificación agrícola implica un cambio
paradigmático en cómo los agricultores asignan sus recursos. El modelo
ofrece una herramienta analítica que trasciende la simple intuición y
las restricciones lineales, al emplear datos históricos sobre
rendimientos, precios y costos para maximizar la rentabilidad esperada
ajustada por riesgo. Al identificar combinaciones de cultivos con baja
correlación, la MPT reduce la volatilidad del portafolio, proporcionando
estabilidad en los ingresos y permitiendo ajustes dinámicos basados en
las condiciones del mercado y las variaciones climáticas.

La Teoría Moderna de Portafolios ofrece una metodología superior a la
programación lineal y al análisis intuitivo de los agricultores. Al
considerar tanto la rentabilidad esperada como el riesgo asociado, e
integrar la diversificación mediante correlaciones entre cultivos, la
MPT mejora significativamente la toma de decisiones en la gestión
agroforestal.




\end{document}
